\documentclass{article}
% if you need to pass options to natbib, use, e.g.:
\PassOptionsToPackage{sort&compress}{natbib}
% before loading neurips_2025

\newcommand{\TODO}[1]{\textbf{\color{red}[TODO: #1]}}
\newcommand{\xun}[1]{{\color{blue}\textbf{Xun:} #1}}
\newcommand{\xuntext}[1]{{\color{blue}#1}}
\newcommand{\zhengqi}[1]{{\color{purple}\textbf{Zhengqi:} #1}}
\newcommand{\guande}[1]{{\color{magenta}\textbf{Guande:} #1}}

\newcommand{\KVSet}{\mathbf{KV}}
\newcommand{\ModelOuput}{\mathbf{X}_{\theta}}
\newcommand{\KVOutput}{\mathbf{kv}}

\usepackage{graphicx}
\usepackage{amsmath, amssymb}
\usepackage{algorithm}
% \usepackage{algpseudocode}
\usepackage{algpseudocode}
\usepackage{multirow}
\usepackage[table]{xcolor}
\usepackage{tablefootnote}
\usepackage{subcaption}

\usepackage{wrapfig}
\usepackage{lipsum} % for filler text
\usepackage{booktabs} % for better table styling
\usepackage{adjustbox}

\definecolor{catgray}{gray}{0.92}
\tracingpages=1

% ready for submission
% \usepackage[sort&compress]{natbib}
% \usepackage[nonatbib]{neurips_2025}


% to compile a preprint version, e.g., for submission to arXiv, add add the
% [preprint] option:
%     \usepackage[preprint]{neurips_2025}


% to compile a camera-ready version, add the [final] option, e.g.:
\usepackage[final]{neurips_2025}


% to avoid loading the natbib package, add option nonatbib:
%    \usepackage[nonatbib]{neurips_2025}


\usepackage[utf8]{inputenc} % allow utf-8 input
\usepackage[T1]{fontenc}    % use 8-bit T1 fonts
\usepackage{hyperref}       % hyperlinks
\usepackage{url}            % simple URL typesetting
\usepackage{booktabs}       % professional-quality tables
\usepackage{amsfonts}       % blackboard math symbols
\usepackage{nicefrac}       % compact symbols for 1/2, etc.
\usepackage{microtype}      % microtypography
\usepackage{xcolor}         % colors
\usepackage{multirow}


% \title{Self Forcing: Bridging Training and Inference in Autoregressive Video Diffusion}
\title{Self Forcing: Bridging the Train-Test Gap in \\ Autoregressive Video Diffusion}
% \title{Eliminating Exposure Bias in \\ Autoregressive Video Diffusion via Self Forcing}
% \title{Self Forcing: Fixing Exposure Bias in Autoregressive Video Diffusion via Holistic Post-Training}
% \title{Self Forcing: Autoregressive Diffusion Post-Training Enables High-Quality Real-Time Video Generation}
% \title{Self Forcing: High-quality Real-Time Video Generation with Autoregressive Diffusion Post-Training}

% \title{Self Forcing: Bridging the Train-Test Gap for Autoregressive Diffusion Video Post-Training}
% \title{Self Forcing: Enabling Real-Time Video Generation with Autoregressive Diffusion Post-Training}

% \title{Self Forcing: Addressing Exposure Bias \\ in Autoregressive Video Diffusion via \\ Holistic Post-Training}
% \title{Self Forcing: Eliminating Exposure Bias in Autoregressive Video Diffusion}

% The \author macro works with any number of authors. There are two commands
% used to separate the names and addresses of multiple authors: \And and \AND.
%
% Using \And between authors leaves it to LaTeX to determine where to break the
% lines. Using \AND forces a line break at that point. So, if LaTeX puts 3 of 4
% authors names on the first line, and the last on the second line, try using
% \AND instead of \And before the third author name.


% \author{%
%   Xun Huang\\
%   Adobe Research \\
%   \texttt{xuhuang@adobe.com} \\
%   \And
%   Zhengqi Li \\
%   Adobe Research  \\
%   \texttt{zhengli@adobe.com} \\
%   \And
%   Guande He \\
%   The University of Texas at Austin \\
%   \texttt{guande.he@utexas.edu} \\
%   \AND
%   Mingyuan Zhou \\
%   The University of Texas at Austin \\
%   \texttt{mingyuan.zhou@mccombs.utexas.edu} \\
%   \And
%   Eli Shechtman \\
%   Adobe Research \\
%   \texttt{elishe@adobe.com}\\
% }
\author{%
Xun Huang$^1$ \quad Zhengqi Li$^1$ \quad Guande He$^2$ \quad Mingyuan Zhou$^2$ \quad Eli Shechtman$^1$ \\
$^1$Adobe Research \quad $^2$The University of Texas at Austin\\\\
% \texttt{xuhuang@adobe.com} \\\\
\centerline{\url{https://self-forcing.github.io/}}
}

\begin{document}


\maketitle


\vspace{-1em}
\begin{abstract}
  We introduce Self Forcing, a novel training paradigm for autoregressive video diffusion models. It addresses the longstanding issue of exposure bias, where models trained on ground-truth context must generate sequences conditioned on their own imperfect outputs during inference. Unlike prior methods that denoise future frames based on ground-truth context frames, Self Forcing conditions each frame’s generation on previously self-generated outputs by performing autoregressive rollout with key-value (KV) caching during training. This strategy enables supervision through a holistic loss at the video level that directly evaluates the quality of the entire generated sequence, rather than relying solely on traditional frame-wise objectives. To ensure training efficiency, we employ a few-step diffusion model along with a stochastic gradient truncation strategy, effectively balancing computational cost and performance. We further introduce a rolling KV cache mechanism that enables efficient autoregressive video extrapolation. Extensive experiments demonstrate that our approach achieves real-time streaming video generation with sub-second latency on a single GPU, while matching or even surpassing the generation quality of significantly slower and non-causal diffusion models.  
\end{abstract}


\section{Introduction}
Recent years have witnessed tremendous progress in video synthesis, with state-of-the-art systems now capable of generating remarkably realistic content with complex temporal dynamics~\cite{videoworldsimulators2024}.
However, these results are typically achieved with diffusion transformers~(DiT)~\cite{peebles2022scalable,wang2025wan} that denoise all frames simultaneously using bidirectional attention. This design allows the future to affect the past and requires generating the entire video at once, fundamentally limiting their applicability to real-time streaming applications where future information is unknown when generating the current frame.

In contrast, autoregressive~(AR) models~\cite{yan2021videogpt,ge2022long,yu2023language,hong2022cogvideo, kondratyuk2024videopoet} generate videos sequentially, a paradigm that naturally aligns with the causal structure of temporal media.
This approach not only significantly reduces the viewing latency of generated videos but also unlocks numerous applications, including real-time interactive content creation~\cite{chen2024streaming,liang2025looking}, game simulation~\cite{oasis2024,yu2025gamefactory,parkerholder2024genie2, valevski2025diffusion}, and robotics learning~\cite{yang2023learning, yu2024learning, li2025unified}.
However, AR models often struggle to match the visual fidelity achieved by state-of-the-art video diffusion models due to their reliance on lossy vector quantization techniques~\cite{van2017neural}.

To combine the best of both worlds, two recent techniques have emerged to equip video diffusion models with AR generation capabilities: Teacher Forcing (TF)~\cite{hu2024acdit,jin2024pyramidal,gao2024ca2,zhang2025test} and Diffusion Forcing (DF)~\cite{chen2024diffusion,yin2025causvid,song2025history,gu2025long,magi1,chen2025skyreels}.
Teacher Forcing, a well-established paradigm in sequence modeling, trains the model to predict the next token conditioned on ground-truth tokens.
When applied to video diffusion, TF involves denoising each frame using clean, ground-truth context frames (Figure~\ref{fig:main} (a)), a strategy commonly referred to as next-frame prediction.
In contrast, Diffusion Forcing trains the model on videos with noise levels independently sampled for each frame, denoising each frame based on noisy context frames~(Figure~\ref{fig:main}~(b)).
This ensures the autoregressive inference scenario, where context frames are clean and the current frame is noisy, is covered by the training distribution.

However, models trained with TF or DF often suffer from error accumulation during autoregressive generation,
leading to degraded video quality over time~\cite{yin2025causvid,zhang2025packing,wang2025error}. This issue is more broadly known as \textit{exposure bias}~\cite{schmidt2019generalization,ning2024elucidating}, where a model is trained exclusively on ground-truth context but must rely on its own imperfect predictions at inference time, resulting in a distributional mismatch that compounds errors as generation progresses. While some approaches attempt to mitigate this issue in video diffusion models by incorporating noisy context frames during inference~\cite{chen2024diffusion,oasis2024,zhang2025packing},
such design sacrifices temporal consistency, complicates the KV-cache design, increases generation latency, and does not fundamentally resolve the exposure bias problem.

\begin{figure*}
  \centering
  % \includegraphics[trim=45pt 10pt 80pt 15pt,width=\textwidth]{figures/overview.pdf}
  \includegraphics[trim=55pt 0pt 105pt 0pt,width=\textwidth]{figures/overview.pdf}
  \caption{
    \textbf{Training paradigms for AR video diffusion models.}
    (a) In Teacher Forcing, the model is trained to denoise each frame conditioned on the preceding clean, ground-truth context frames.
    (b) In Diffusion Forcing, the model is trained to denoise each frame conditioned on the preceding context frames with varying noise levels.
    Both (a) and (b) generate outputs that do not belong to the distribution the model generates during inference.
    (c) Our Self Forcing approach performs autoregressive self-rollout \textit{during training}, denoising the next frame based on previous context frames generated by itself. A distribution-matching loss~(\textit{e.g.}, SiD, DMD, GAN) is computed on the final output video to align the distribution of generated videos with that of real videos.
    Our training paradigm closely mirrors the inference process, thereby bridging the train-test distribution gap.
  }
  \vspace{-1.5em}
  \label{fig:main}
\end{figure*}


In this work, we propose \textit{Self Forcing}~(SF), a novel algorithm addressing exposure bias in autoregressive video generation.
Inspired by early RNN-era sequence modeling techniques~\cite{lamb2016professor,ranzato2016sequence,yu2017seqgan}, our approach bridges the train-test distribution gap by explicitly unrolling autoregressive generation during training, generating each frame conditioned on previously self-generated frames rather than ground-truth ones.
This enables supervision with holistic distribution-matching losses~\cite{yin2024one,yin2024improved,goodfellow2014generative} applied to complete generated video sequences.
By forcing the model to encounter and learn from its own prediction errors, Self Forcing effectively mitigates exposure bias and reduces \mbox{error accumulation.}

While Self Forcing may seem computationally prohibitive due to its sequential nature preventing parallel training, we demonstrate that it can be efficiently implemented as an algorithm in the post-training stage where the model does not require a large number of gradient updates to converge.
By employing a few-step diffusion backbone and a carefully designed gradient truncation strategy, Self Forcing is surprisingly more efficient than alternative parallel strategies, achieving superior performance within the same wall-clock training time.
Additionally, we introduce a rolling KV cache mechanism that enhances the efficiency of video extrapolation.

Extensive experiments demonstrate that our model enables real-time video generation at 17 FPS with sub-second latency on a single H100 GPU, while achieving competitive or superior generation quality compared to recent slow bidirectional and autoregressive video diffusion models.
These advances open the door to genuinely interactive video generation use cases—live streaming, gaming, and world simulation—where latency budgets are measured in milliseconds rather than minutes.


\section{Related Work}
% \section{Related Work and Preliminaries}
% GANs for Video Generation
\noindent\textbf{GANs for Video Generation.}
Early video generation approaches relied primarily on generative adversarial networks~(GANs)~\cite{goodfellow2014generative}, either using convolutional networks to generate entire videos in parallel~\cite{vondrick2016generating,saito2017temporal,brooks2022generating} or employing recurrent architectures to produce frames sequentially~\cite{villegas2017decomposing,denton2017unsupervised,tulyakov2018mocogan,liu2021infinite, li2022infinitenature}.
Recently, GANs have also been applied to distill video diffusion models~\cite{zhang2024sf,lin2025diffusion,mao2025osv,wu2025snapgen}.
Since the generator in GANs follows the same process during training and inference, it inherently avoids exposure bias. Our work draws inspiration from this fundamental GAN principle by directly optimizing the alignment between the generator's output distribution and the target distribution.
% \noindent\textbf{GANs for Video Generation.}
% Early video generation approaches relied primarily on GANs~\cite{goodfellow2014generative}, either using convolutional architectures to generate videos in parallel~\cite{vondrick2016generating,saito2017temporal,brooks2022generating} or employing recurrent networks for sequential frame production~\cite{villegas2017decomposing,denton2017unsupervised,tulyakov2018mocogan,liu2021infinite,li2022infinitenature}. Recently, GANs have been repurposed to distill video diffusion models~\cite{zhang2024sf,lin2025diffusion,mao2025osv,wu2025snapgen}. A key strength of GANs lies in their training-inference alignment—the generator follows identical processes in both phases, inherently preventing distribution mismatch. Our work draws inspiration from this fundamental GAN principle by directly optimizing alignment between the generator's output distribution and the target distribution during training.

% Diffusion/AR for Video Generation
% \vspace{1em}
\noindent\textbf{Autoregressive/Diffusion Models for Video Generation.}
Modern video generation models have largely shifted toward diffusion or autoregressive models due to their stronger scaling abilities. Video diffusion models typically adopt bidirectional attention mechanisms to simultaneously denoise all video frames~\cite{ho2022video,blattmann2023align,Ho2022ImagenVH,videoworldsimulators2024,polyak2024movie,yang2025cogvideox,hacohen2024ltx,kong2024hunyuanvideo,blattmann2023stable,villegas2023phenaki,deng2024nova,gupta2024photorealistic,wang2025wan}. Autoregressive models, in contrast, are trained with next-token prediction objectives and generate spatiotemporal tokens sequentially at inference time~\cite{weissenborn2020scaling, kondratyuk2024videopoet,yan2021videogpt,wang2024loong,Bruce2024GenieGI,ren2025next}.
% ho2022video, ,blattmann2023stable,
% They usually rely on vector quantization~(VQ) techniques~\cite{van2017neural} to discretize visual tokens and apply cross-entroy loss.
% cross-entropy loss.

% AR Diffusion hybrid
% \vspace{1em}
\noindent\textbf{Autoregressive-Diffusion Hybrid Models.}
Very recently, hybrid models integrating autoregressive and diffusion frameworks have emerged as a promising direction in generative modeling of videos~\cite{weng2024art,liu2024redefining,chen2024diffusion,guo2025long,yin2025causvid,hu2024acdit,jin2024pyramidal,gu2025long,gao2024ca2,li2025arlon,liu2024mardini,zhang2025generative,zhang2025test} as well as other sequence domains \cite{li2024autoregressive,wu2023ar,deng2024causal,arriola2025block,liu2024autoregressive,zhou2025transfusion,mo2025xfusionintroducingnewmodality}.
% These approaches autoregressively decompose the distribution of video frames $p(\{\mathbf{x}_i\})$ into product of conditionals $\prod_{i=1}^{N} p(\mathbf{x}_i | \mathbf{x}_{<i})$, where each conditional $p(\mathbf{x}_i | \mathbf{x}{<i})$ is modeled with diffusion.
They typically rely on a long, iterative prediction chain~(both temporally autoregressive and spatially iterative denoising), which could lead to significant error accumulation.
% , especially since the model is trained only on ground-truth but not on generated data.
% where each prediction is conditioned on previous predictions. This leads to
% , commonly known as the problem of exposure bias.
Our work addresses this issue
% in autoregressive diffusion models
by training the model conditioned on its own predictions and teaching it to correct its own mistakes.

%  Rolling diffusion style
% \vspace{1em}
\noindent\textbf{Rolling Diffusion and Variants.}
Another line of work~\cite{ruhe2024rolling,kim2025fifo,xie2024progressive,zhang2025packing,magi1,sun2025ar} trains video diffusion models with a progressive noise schedule, where the noise level gradually increases from earlier to later frames. While these methods support sequential long video generation with less accumulated errors and are sometimes also referred to as autoregressive, they do not strictly follow the autoregressive chain rule decomposition. Consequently, they would exhibit significant latency in interactive applications, as future frames are partially pre-generated before the current frame is presented to the user. This premature commitment restricts the impact of real-time user-injected controls, resulting in limited responsiveness in immediately subsequent frames.

% \vspace{1em}
\noindent\textbf{CausVid.}
Our work is most closely related to CausVid~\cite{yin2025causvid}, which trains few-step autoregressive diffusion models using the DF scheme and distribution matching distillation (DMD).  However, CausVid suffers from a critical flaw that its training outputs (generated via DF) do not come from the distribution the model produces at inference time, therefore the DMD loss is matching the wrong distribution. We pinpoint this issue and propose a solution that matches the true model distribution.

\section{Self Forcing: Briding Train-Test Gap via Holistic Post-Training}
We first provide a formal definition of autoregressive video diffusion models and describe standard training approaches in Section~\ref{sec:preliminaries}. In Section~\ref{sec:rollout}, we introduce the main part of our Self Forcing training algorithm and describe how it can be efficiently implemented with a few-step diffusion model.
In Section~\ref{sec:loss}, we describe various choices of holistic, video-level distribution-matching training objectives. Finally, we introduce a rolling key-value cache mechanism that enables efficient generation of arbitrarily long videos in Section~\ref{sec:rolling_kv}.

\subsection{Preliminaries: Autoregressive Video Diffusion Models}
\label{sec:preliminaries}
Autoregressive video diffusion model is a hybrid generative model that combines autoregressive chain-rule decomposition with denoising diffusion models for video generation. % Xun: let's keep the first sentence high-level.
Specifically, given a sequence of $N$ video frames \( x^{1:N} = (x^1, x^2, \dots, x^N) \), it factorizes the joint distribution into product of conditionals using the chain rule $p(x^{1:N}) = \prod_{i=1}^{N} p(x^i|x^{<i})$.
Each conditional distribution \( p(x^i|x^{<i}) \) is then modeled using a diffusion process, where a frame is generated by progressively denoising an initial Gaussian noise conditioned on previously generated frames.
This formulation combines the strengths of both autoregressive models and diffusion models for capturing sequential dependencies while enabling high-quality generation of continuous-valued visual signal.
In practice, we can also choose to generate one chunk of frames rather than a single frame at a time~\cite{yin2025causvid,magi1}. For simplicity of notation, however, we continue to denote each chunk as a frame in this section.

% How prior work trains autoregressive video diffusion models?
Most existing autoregressive video diffusion models are trained using frame-wise denoising loss within the paradigm of Teacher Forcing (TF) or Diffusion Forcing (DF).
Specifically, each frame \(x^i\) is corrupted by the forward process \( q_{t^i|0}(x_{t^i}^{i} | x_0^{i}) \) such that \(x_{t^i}^i  = \Psi(x^i, \epsilon^i, t^i) = \alpha_{t^i} x^i + \sigma_{t^i} \epsilon^i \), where \( \alpha_{t^i}, \sigma_{t^i} \) are pre-defined noise schedule within a finite time horizon \( t^i \in [0, 1000] \) and \( \epsilon^i \sim \mathcal{N}(0, I) \) is Gaussian noise. In TF, the timesteps $t^i$ are usually shared across all frames, whereas in DF, they are sampled independently for each frame.
A generative model is learned through the time-reversal of the forward process, where each denoising step can be achieved by predicting the noise $\epsilon^i$ added to each frame with a neural network \(\hat \epsilon_\theta^i := G_\theta(x_{t^i}^i, t^i, c) \) conditioned on the context $c$.
The context consists of clean ground-truth frames \( x^{<i} \) in TF or noisy context frames \( x^{j<i}_{t^j} \) in DF.
% \zhengqi{shall we also mention other moedialty like text condition here?}
% , where \( c \) is the clean context frames \( x^{<i} \) for TF or the noisy context frames \( x^{j<i}_{t^j} \) for DF. \guande{There is a double where phrase}
The model is trained to minimize the frame-wise mean squared error (MSE) between the predicted noise and the true added noise:
\( \mathcal{L}^{\text{DM}}_\theta = \mathbb{E}_{x^i, t^i, \epsilon^i} \left[ w_{t^i} \| \hat{\epsilon}^i_\theta - \epsilon^i \|_2^2 \right] \), where $w_{t^i}$ is a weighting function.

\begin{figure*}
  \centering
  \includegraphics[trim=35pt 5pt 65pt 5pt,width=\textwidth]{figures/attention.pdf}
  \caption{
    \textbf{Attention mask configurations.}  Both Teacher Forcing (a) and Diffusion Forcing (b) train the model on the entire video in parallel, enforcing causal dependencies with custom attention masks. In contrast, our Self-Forcing Training (c) mirrors the autoregressive (AR) inference process with KV caching and does not rely on special attention masks. For illustration purposes, we show a scenario where the video contains 3 frames, and each frame consists of 2 tokens.
  }
  \vspace{-1em}
  \label{fig:attention}
\end{figure*}

We focus on the transformer-based architecture~\cite{peebles2022scalable} of diffusion models with text conditioning~(omitted from equations for clarity) operating in a compressed latent space encoded by a causal 3D variational autoencoder (VAE)~\cite{kingma2013auto}. The autoregressive chain-rule decomposition is implemented via causal attention. Figures~\ref{fig:attention} (a) and (b) illustrate the attention mask configurations of Teacher Forcing and Diffusion Forcing approaches. For Teacher Forcing, we describe an efficient variant that processes all frames in parallel using block sparse attention masks, rather than denoising one frame at each training iteration~\cite{jin2024pyramidal}. Such design has been used in MAR-based~\cite{li2024autoregressive} autoregressive video generation~\cite{zhou2025taming} and concurrently in other autoregressive video diffusion models~\cite{zhang2025generative,zhang2025test}.


\subsection{Autoregressive Diffusion Post-Training via Self-Rollout}
\label{sec:rollout}
The core idea of Self Forcing is to generate videos through autoregressive self-rollout during training following the inference-time recipe. Specifically, we sample a batch of videos $\{x^{1:N}_{\theta}\}\sim p_\theta(x^{1:N}) = \prod_{i=1}^{N} p_\theta(x^i|x^{<i})$ where each frame \( x^i \) is generated by performing iterative denoising conditioned on self-generated outputs, including both clean context frames in the past and noisy frames at the current time step.
Unlike most previous autoregressive models that only utilize KV caching during inference, our Self Forcing method innovatively employs KV caching during training, as shown in Figure~\ref{fig:attention}~(c).

Nevertheless, implementing Self Forcing with standard many-step diffusion models would be computationally prohibitive, as it requires unrolling and backpropagation through long denoising chains.
Therefore, we choose to use a few-step diffusion model $G_\theta$ to approximate each conditional distribution $p_\theta(x^i|x^{<i})$ in the autoregressive factorization. Consider $\{t_0=0,t_1, \dots, t_{T}=1000\}$ a subsequence of timesteps $[0, ..., 1000]$,
at each denoising step $t_j$ and frame index $i$, the model denoises an intermediate noisy frame \( x^i_{t_j} \) conditioned on previous clean frames \( x^{<i} \). It then injects Gaussian noise with a lower noise level into the denoised frame through the forward process $\Psi$ to obtain the noisy frame \( x^i_{t_{j - 1}} \) as the input to the next denoising step, following the standard practice in few-step diffusion models~\cite{song2023consistency,yin2024improved}.
The model distribution $p_\theta(x^i|x^{<i})$ is implicitly defined as $ f_{\theta,t_1} \circ f_{\theta,t_2} \circ ... \circ f_{\theta,t_T} (x^i_{t_T})$, where $f_{\theta,t_j}(x_{t_j}^i)=\Psi(G_\theta(x^i_{t_j}, t_{j},  x^{<i}), \epsilon_{t_{j-1}}, t_{j-1})$, \mbox{and $ x^i_{t_{T}} \sim \mathcal{N}(0, I)$}.
\begin{figure}[t!]
\begin{minipage}[t]{0.48\textwidth}
  \begin{algorithm}[H]
    \caption{Self Forcing Training}
    \small
    \begin{algorithmic}[1]
      \Require Denoise timesteps $\{t_1, \dots, t_T\}$
      \Require Number of video frames $N$
      \Require AR diffusion model $G_\theta$ (returns KV embeddings via $G_\theta^{\mathrm{KV}}$)
      \Loop
        \State Initialize model output $\ModelOuput \gets []$
        \State Initialize KV cache $\KVSet \gets []$
        \State Sample $s \sim \text{Uniform}(1, 2, \ldots, T)$
        \For{$i = 1, \dots, N$}
          \State Initialize $x^i_{t_T} \sim \mathcal{N}(0, I)$
          \For{$j = T, \dots, s$}
            \If{$j = s$}
              \State Enable gradient computation
              \State Set $\hat{x}^i_{0} \gets G_\theta(x^i_{t_{j}}; t_j, \KVSet)$
              \State $\ModelOuput{\texttt{.append}}(\hat x^i_{0})$
              \State Disable gradient computation
              \State Cache $\KVOutput^i \gets G_\theta^\text{KV}(\hat{x}^i_{0}; 0, \KVSet)$
              \State $\KVSet{\texttt{.append}}(\KVOutput^i)$
            \Else
              \State Disable gradient computation
              \State Set $\hat{x}^i_{0} \gets G_\theta(x^i_{t_j}; t_j, \KVSet)$
              \State Sample $\epsilon \sim \mathcal{N}(0, I)$
              \State Set $x^i_{t_{j-1}} \gets \Psi(\hat{x}^i_0, \epsilon, t_{j-1})$
            \EndIf
          \EndFor
        \EndFor
        \State Update $\theta$ via distribution matching loss
      \EndLoop
    \end{algorithmic}
    \label{alg:self_forcing}
  \end{algorithm}
\end{minipage}
\hfill
\begin{minipage}[t]{0.48\textwidth}
  \begin{algorithm}[H]
    \caption{Autoregressive Diffusion Inference with Rolling KV Cache}
    \small
    \begin{algorithmic}[1]
      \Require KV cache of size $L$ frames
      \Require Denoise timesteps $\{t_1, \dots, t_T\}$
      \Require Number of generated frames $M$
      \Require AR diffusion model $G_\theta$ (returns KV embeddings via $G_\theta^\text{KV}$)
      \State Initialize model output $\ModelOuput \gets []$
      \State Initialize KV cache $\KVSet \gets []$
      \For{$i = 1, \dots, M$}
        \State Initialize $x^i_{t_T} \sim \mathcal{N}(0, I)$
        \For{$j = T, \dots, 1$}
          \State Set $\hat{x}^i_{0} \gets G_\theta(x^i_{t_j}; t_j, \KVSet)$
          \If{$j = 1$}
            \State $\ModelOuput{\texttt{.append}}(\hat x^i_{0})$
            \State Cache $\KVOutput^i \gets G_\theta^\text{KV}(\hat{x}^i_{0}; 0, \KVSet)$
            \If{$|\KVSet| = L$}
              \State $\KVSet.\mathrm{pop}(0)$ \Comment{Cache eviction}
            \EndIf
            \State $\KVSet{\texttt{.append}}(\KVOutput^i)$
          \Else
            \State Sample $\epsilon \sim \mathcal{N}(0, I)$
            \State Set $x^i_{t_{j-1}} \gets \Psi(\hat{x}^i_0, \epsilon, t_{j-1})$
          \EndIf
        \EndFor
      \EndFor
      \State \Return $\ModelOuput$
    \end{algorithmic}
    \label{alg:inference}
  \end{algorithm}
\end{minipage}
\vspace{-1em}
\end{figure}



Even with few-step models, however, naively backpropagating through the entire autoregressive diffusion process would still lead to excessive memory consumption.
To address this challenge, we propose a gradient truncation strategy that limits the backpropagation to only the final denoising step of each frame. Moreover, instead of always using $T$ denoising steps~(as in inference time), we randomly sample a denoising step $s$ from $[1, T]$ for each sample sequence at each training iteration, and use the denoised output of the $s$-th step as the final output.
This stochastic sampling approach ensures all intermediate denoising steps receive supervision signals.
We additionally detach the gradients of the previous frames from the current frame during training by restricting gradient flow into KV cache embeddings. For a complete description of the training process, see Algorithm~\ref{alg:self_forcing}.


\subsection{Holistic Distribution Matching Loss}
\label{sec:loss}
Autoregressive self-rollout generates samples directly from the inference-time model distribution, enabling us to apply holistic, video-level losses that align the distribution of generated videos \( p_\theta(x^{1:N}) \) with that of real videos \( p_{\text{data}}(x^{1:N}) \).
To leverage pre-trained diffusion models and enhance training stability~\cite{jenni2019stabilizing}, we inject noise to both distributions and match \( p_{\theta, t}(x^{1:N}_t) \) and \( p_{\text{data}, t}(x^{1:N}_t) \), where each represents the respective distribution after applying the forward diffusion process: \( p_{\cdot, t}(x^{1:N}_t) = \int q_{t|0}(x^{1:N}_t | x^{1:N}) p_{\cdot}(x^{1:N}) \mathrm{d} x^{1:N} \).
Our framework is generally applicable to various divergence measures and distribution matching frameworks, and we consider three approaches in this paper:
\begin{itemize}


  \item \textbf{Distribution Matching Distillation (DMD)}~\cite{yin2024one,yin2024improved}: This approach minimizes the reverse Kullback-Leibler divergence $\mathbb{E}_{t} [ D_\text{KL}(p_{\theta, t} \| p_{\text{data}, t}) ]$ by leveraging the score difference between distributions to guide gradient updates.
  \item \textbf{Score Identity Distillation (SiD)}~\cite{zhou2024score,zhou2025adversarial}: This method performs distribution matching via Fisher divergence $\mathbb E_{t, p_{\theta, t}}[\| \nabla \log p_{\theta, t}\ - \nabla \log p_{\text{data}, t}\|^2]$.

  \item \textbf{Generative Adversarial Networks (GANs)}~\cite{goodfellow2014generative}: It approximately minimizes the Jensen-Shannon divergence through a minimax game between the generator (our autoregressive diffusion model) and a discriminator that distinguishes between real and generated videos.
\end{itemize}

Importantly, our training objective matches the \textit{holistic} distribution of the entire video sequence to the data distribution $D(p_{\text{data}}(x^{1:N}) \| p_\theta(x^{1:N}))$. In contrast, TF/DF can be understood as performing \textit{frame-wise} distribution matching: $\mathbb{E}_{\{x^{<i}\}\sim p_{\text{data}}} D_{KL}(p_{\text{data}}(x^{i}|x^{<i}) \| p_\theta(x^{i}|x^{<i}))$\footnote{With a specific weighting per noise level~\cite{kingma2021variational,song2021maximum}, denoising loss approximates the maximum likelihood objective, equivalent to minimizing the KL divergence between per-frame data and model distributions.}, where DF additionally samples context frames from a noise-corrupted data distribution $\{x^{<i}\}\sim\tilde{p}_{\text{data}}$.
Our formulation fundamentally transforms the training dynamics—context frames $\{x^{<i}\}$ are sampled from the model's own distribution $p_\theta$ rather than from the data distribution (clean or noisy). This alignment between training and inference distributions effectively addresses exposure bias and forces the model to learn from its own imperfections, thereby developing robustness to error accumulation.

While all three objectives have been used in the context of timestep distillation of diffusion models, our primary motivation differs fundamentally from distillation: we aim to enhance the \textit{quality} of autoregressive video generation by addressing exposure bias via distribution matching, rather than merely accelerating sampling.
This distinction makes other popular distillation methods~\cite{song2023consistency} inapplicable to our framework as they only focus on timestep reduction without directly aligning the generator output distribution.
Although CausVid~\cite{yin2025causvid} similarly employs DMD to match the distribution of generated videos, the distribution it optimizes during training (using Diffusion Forcing outputs) deviates from the actual inference-time distribution, significantly undermining its effectiveness.

\subsection{Long Video Generation with Rolling KV Cache}
\label{sec:rolling_kv}

A key advantage of autoregressive models over standard video diffusion models is their extrapolative ability, in principle allowing the generation of infinitely long videos via sliding-window inference.
While bidirectional attention models trained with Diffusion Forcing~\cite{song2025history,chen2025skyreels} can also generate videos autoregressively, they do not support KV caching, requiring complete recomputation of attention matrices for each new frame. This leads to excessive computational complexity of $O(TL^2)$ (where $T$ represents the number of denoising steps and $L$ is thewindow size), as shown in Figure~\ref{fig:kv}~(a).

Models with causal attention, on the other hand, can leverage KV caching to improve efficiency.
However, existing implementations~\cite{yin2025causvid,magi1} require recomputing KV cache for overlapping frames between consecutive sliding windows, as illustrated in Figure~\ref{fig:kv}~(b).
This leads to $O(L^2+TL)$ complexity when employing dense sliding windows.
%  (incrementing one frame at a time)
As a result, prior implementations adopt larger strides with minimal overlap to reduce computational costs, which compromises temporal consistency since frame at the beginning of each window relies on a significantly reduced historical context.


Inspired by research in large language models~\cite{streamingllm}, we propose a rolling KV cache mechanism for autoregressive diffusion models that allows infinitely long video generation without any need of recomputing the KV cache. As illustrated in Figure~\ref{fig:kv}~(c), we maintain a fixed-size KV cache that stores the KV embeddings of tokens in the most recent $L$ frames. When generating a new frame, we first check if the KV cache is full. If it is, we remove the oldest KV cache entry before adding the new one. This approach enables endless frame generation with a time complexity of $O(TL)$, while still maintaining a sufficient context length when generating each new frame. Algorithm~\ref{alg:inference} provides a detailed description of our autoregressive long video generation algorithm with rolling KV cache.


\begin{figure*}[t]
  \centering
  \includegraphics[trim=15pt 10pt 15pt 10pt,width=\textwidth]{figures/rolling_kv.pdf}
  \caption{
    \textbf{Efficiency comparisons for video extrapolation.}
    When performing video extrapolation through sliding window inference, (a) bidirectional diffusion models trained with TF/DF~\cite{song2025history,chen2025skyreels} do not support KV cache. (b) Prior causal diffusion models~\cite{yin2025causvid,magi1} require re-computing KV when shifting the window. (c) Our method does not recompute KV and enables \mbox{more efficient extrapolation.}
  }
  \label{fig:kv}
  \vspace{-1em}
\end{figure*}

However, naive implementation of this mechanism leads to severe flickering artifacts due to distribution mismatch. Specifically, the first latent frame has different statistical properties than other frames: it only encodes the first image without performing temporal compression. The model, having always seen the first frame as the image latent during training, fails to generalize when the image latent is no longer visible in the rolling KV cache scenario.
% While we experimented with strategies similar to StreamingLLM~\cite{streamingllm}, keeping the first frame fixed while rolling KV cache of other frames, these proved ineffective for video generation. 
Our solution is straightforward but effective: during training, we restrict the attention window so the model cannot attend to the first chunk when denoising the final chunk, thereby simulating the conditions encountered during long video generation.



\section{Experiments}
\vspace{-0.5em}
\label{sec:experiments}

\noindent\textbf{Implementation.}
We implement Self Forcing with Wan2.1-T2V-1.3B~\cite{wang2025wan}, a Flow Matching~\cite{lipmanflow} based model that generates 5s videos at 16 FPS with a resolution of $832\times480$.
Following CausVid's initialization protocol~\cite{yin2025causvid}, we first finetune the base model with causal attention masking on 16k ODE solution pairs sampled from the base model.
For both ODE initialization and Self Forcing training, we sample text prompts from a filtered and LLM-extended version of VidProM~\cite{wangvidprom}.
We use 4-step diffusion and implement both frame-wise and chunk-wise autoregressive variants, with the latter generating a chunk of 3 latent frames at a time.
We adopt the R3GAN~\cite{huang2024gan} objective, which consists of relativistic pairing GAN loss~\cite{jolicoeurrelativistic} with R1 + R2 regularization~\cite{mescheder2018training}.
We use the 14B base model to generate 70k videos as the dataset for training GANs~\cite{sauer2024fast} and fine-tuning many-step TF/DF AR diffusion baselines. Notably, DMD/SiD implementations of our algorithm remain data-free, capable of converting a pre-trained video diffusion model into an autoregressive model without any video training data.
Additional implementation details are provided in Appendix~\ref{app:implementation}.

\noindent\textbf{Evaluation metrics.}
We adopt VBench~\cite{huang2023vbench} and user preference study to evaluate both visual quality and semantic alignment.
We also rigorously evaluate the efficiency of our method for real-time applications.
While some recent works claim ``real-time'' video generation abilities~\cite{hacohen2024ltx,zhao2025real} based solely on \textit{throughput}, we argue that true real-time performance requires both sufficient throughput (exceeding video playback rate) and lower \textit{latency} than the perceptual threshold which could be application-dependent~\cite{li4654242latency}.
We therefore evaluate both throughput and first-frame latency to provide a comprehensive assessment of real-time capabilities, with all speed tests conducted on a single NVIDIA H100 GPU.


\begin{wrapfigure}{r}{0.48\textwidth}
  \vspace{-1em}
  \centering
  \includegraphics[width=\linewidth]{figures/user_study_result.pdf}
  \caption{\textbf{User preference study.} Self Forcing outperforms all baselines in human preference.}
  \label{fig:user_study_result}
  \vspace{-1em}
\end{wrapfigure}

\noindent\textbf{Comparison with existing baselines.}
We compare our model with relevant open-source video generation models of similar scale. Our comparisons include two diffusion models: Wan2.1-1.3B~\cite{wang2025wan} (our initialization weights) and LTX-Video~\cite{hacohen2024ltx} (known for efficiency). We also compare with several autoregressive models including Pyramid Flow~\cite{jin2024pyramidal}, NOVA~\cite{deng2024nova}, SkyReels-V2~\cite{chen2025skyreels}, MAGI-1~\cite{magi1}, and CausVid~\cite{yin2025causvid} (also initialized from Wan-1.3B).

As shown in Table~\ref{tab:main}, our chunk-wise autoregressive model achieves the highest VBench scores across all compared models while simultaneously delivering real-time throughput ($17.0$ FPS) with sub-second latency, low enough for certain real-time applications such as live video streaming~\cite{bentaleb2025toward}.
Figure~\ref{fig:user_study_result} shows the user study results comparing our chunk-wise Self Forcing model against several important baselines. Our approach is consistently preferred over all alternatives, including the many-step diffusion model Wan2.1 that our model is initialized from.
Our frame-wise variant maintains strong generation quality while providing the lowest latency ($0.45$s), making it particularly suitable for latency-sensitive real-time applications.
% Results in this table are obtained with the DMD loss, but we observe similar performance with SiD and GAN losses and report them in the following.
Results here are obtained using the DMD loss objective; models trained with SiD and GAN objectives achieve comparable performance as detailed in our ablation studies.
As shown in Figure~\ref{fig:comp_5s}, CausVid suffers from the error accumulation problem that causes the saturation to increase over time. Our approach obtains slightly better visual quality than Wan2.1/SkyReels-V2, while being around \textbf{150x} faster in latency.
More example videos are provided in the project website (\url{https://self-forcing.github.io/}).
% \noindent\textbf{Comparison with existing models.}
% We compare our approach against state-of-the-art open-source video generation models of comparable scale. Our comparisons include the Wan-1.3B base model~\cite{wang2025wan} (our initialization point), LTX-Video~\cite{hacohen2024ltx} (known for efficiency), and several autoregressive architectures including Pyramid Flow~\cite{jin2024pyramidal}, NOVA~\cite{deng2024nova}, SkyReels-V2~\cite{chen2025skyreels}, MAGI-1~\cite{magi1}, and CausVid~\cite{yin2025causvid} (also Wan-1.3B initialized).
% As shown in Table~\ref{tab:main}, our chunk-wise autoregressive model achieves the highest VBench scores across all compared models while simultaneously delivering real-time throughput (17.4 FPS) with sub-second initial latency. Our frame-wise variant maintains strong generation quality while providing the lowest latency among all tested models (0.44s), making it particularly suitable for interactive streaming applications requiring immediate visual feedback.
% and other autoregressive models including CausVid~\cite{yin2025causvid} and SkyReels-V2~\cite{chen2025skyreels}. We also include the LTX-Video model as a representative of closed-source models. As shown in Tab.~\ref{tab:main}, our chunk-wise autoregressive model achieves better performance than the base Wan2.1 diffusion model as well as other autoregressive models. Our frame-wise autoregressive model achieves significantly lower latency and better quality than all previous autoregressive models.

\begin{table}[t]
  \small
  \setlength{\tabcolsep}{3.5pt} % tighten columns a bit more to accommodate new column
  \caption{
    \textbf{Comparison with relevant baselines.} We compare Self Forcing with representative open-source video generation models of similar parameter sizes and resolutions.
    % Our chunk-wise AR model achieves better performance than the base Wan2.1 diffusion model as well as other AR models. Our frame-wise AR model achieves significantly lower latency and better quality than all previous AR models.
  }
  \vspace{0.5em}
  \label{tab:main}
  \centering
  \begin{tabular}{lcccccccc}
      \toprule
      \multirow{2}{*}{Model} & \multirow{2}{*}{\#Params} & \multirow{2}{*}{Resolution} & \multirow{2}{*}{Throughput} & \multirow{2}{*}{Latency} &
      \multicolumn{3}{c}{Evaluation scores $\uparrow$}\\
    \cmidrule(lr){6-8}
     & & & (FPS) $\uparrow$ & (s)  & Total & Quality & Semantic \\
     & & &                  &                  & Score & Score & Score \\
    \midrule
    % \rowcolor{catgray}
    % \multicolumn{8}{l}{\textit{Closed-source models}}\\
    % Sora                      & -- & -- & --   & -- & 84.28 & 85.51 & 79.35 \\
    % \midrule
    \rowcolor{catgray}
    \multicolumn{8}{l}{\textit{Diffusion models}}\\
    LTX-Video~\cite{hacohen2024ltx}      & 1.9B & $768{\times}512$ & 8.98          & 13.5          & 80.00 & 82.30 & 70.79 \\
    Wan2.1~\cite{wang2025wan}            & 1.3B & $832{\times}480$ & 0.78          & 103           & 84.26 & \textbf{85.30} & 80.09 \\
    % HunyuanVideo                       & 13B  & -- & --   & 48 & 83.24 & 85.09 & 75.82 \\
    \midrule
    \rowcolor{catgray}
    \multicolumn{8}{l}{\textit{Chunk-wise autoregressive models}}\\
    SkyReels-V2~\cite{chen2025skyreels}  & 1.3B & $960{\times}540$ & 0.49          & 112           & 82.67 & 84.70 & 74.53 \\
    MAGI-1~\cite{magi1}                  & 4.5B & $832{\times}480$ & 0.19          & 282           & 79.18 & 82.04 & 67.74 \\
    CausVid~\cite{yin2025causvid}$^{*}$  & 1.3B & $832{\times}480$ & \textbf{17.0} & 0.69          & 81.20 & 84.05 & 69.80 \\
    \midrule
    Self Forcing (Ours, chunk-wise)      & 1.3B & $832{\times}480$ & \textbf{17.0} & 0.69          & \textbf{84.31} & 85.07 & \textbf{81.28}    \\  % 17.0, 0.69
    \midrule
    \rowcolor{catgray}
    \multicolumn{8}{l}{\textit{Autoregressive models}$^{\dagger}$}\\
    NOVA~\cite{deng2024nova}             & 0.6B & $768{\times}480$ & 0.88          & 4.1  & 80.12 & 80.39 & 79.05 \\
    Pyramid Flow~\cite{jin2024pyramidal} & 2B   & $640{\times}384$ & 6.7           & 2.5           & 81.72 & 84.74 & 69.62 \\
    \midrule
    Self Forcing (Ours, frame-wise)      & 1.3B & $832{\times}480$ & 8.9           & \textbf{0.45} & 84.26 & 85.25 & 80.30 \\ % 8.9, 0.45
    \bottomrule
  \end{tabular}\\
  \vspace{2pt}
  {\footnotesize\raggedright$^{*}$ We compare with the official implementation of CausVid that uses the same base model~(Wan-1.3B).\par}
  {\footnotesize\raggedright$^{\dagger}$ The distinction of AR/non-AR applies to the temporal dimension.\par}
  % We do not distinguish different ways of generating tokens within the same frame.
  \vspace{-1em}
\end{table}


\begin{figure}[t!]
  \centering
  \includegraphics[trim=100pt 1pt 1pt 1pt, width=\textwidth]{figures/Self-Forcing-Baseline-2-compressed.pdf}
  \caption{\textbf{Qualitative comparisons.} We visualize videos generated by Self Forcing~(Ours) against those by Wan2.1~\cite{wang2025wan}, SkyReels-V2~\cite{chen2025skyreels}, and CausVid~\cite{yin2025causvid} at three time steps. All models share the same architecture with 1.3B parameters.}
  \label{fig:comp_5s}
  \vspace{-1.5em}
\end{figure}



\noindent\textbf{Ablation Studies.}
We perform controlled comparisons  of Self Forcing with alternative autoregressive diffusion training approaches. We evaluate: (1) AR Diffusion models trained with denoising diffusion loss using either Teacher Forcing or Diffusion Forcing, and (2) few-step AR Diffusion models trained with TF/DF inputs but optimized with distribution matching objectives. The latter configuration with DF and DMD essentially replicates CausVid~\cite{yin2025causvid} within our implementation framework, allowing direct comparison under identical training conditions.
% We initialize the AR Diffusion models with Wan-1.3B weights and the few-step models with the same weights from ODE regression as Self Forcing.
% We report results obtained with Self Forcing using different distribution matching objectives (DMD, SiD, and GAN).

Table~\ref{tab:ablation} demonstrates that Self Forcing performs robustly across various distribution matching objectives (DMD, SiD, and GAN), consistently outperforming all baselines.
While baseline methods exhibit notable quality degradation when shifting from chunk-wise to frame-wise AR due to error accumulation associated with increased AR unrolling steps, usually manifesting as progressive over-saturation or over-sharpening (similar to CausVid in Appendix~\ref{app:qualitative} Fig.~\ref{fig:comp_5s}), Self Forcing maintains consistent performance across both setups, highlighting its effectiveness at addressing exposure bias.

% There are two interesting observations. First,  Second, our results indicate that Teacher Forcing generally surpasses Diffusion Forcing, suggesting it may serve as a more effective strategy for AR diffusion pre-training while Self Forcing can be employed in post-training. \xun{Not sure if this would offend DF people who might be our reviewers...}
% Additional visual comparisons can be found in the supplementary material.
% Table~\ref{tab:ablation} demonstrates that Self Forcing performs robustly across various distribution matching objectives (DMD, SiD, and GAN), consistently outperforming all baseline methods. Two key observations emerge: First, while baseline approaches exhibit notable quality degradation when transitioning from chunk-wise to frame-wise architectures—due to heightened error accumulation manifesting as saturation or sharpness artifacts—Self Forcing maintains stable performance across both setups, underscoring its resilience to exposure bias. Second, our results indicate that Teacher Forcing generally surpasses Diffusion Forcing, suggesting it may serve as a more effective pre-training strategy for autoregressive diffusion models. Additional qualitative comparisons are included in the supplementary material.

% Table~\ref{tab:ablation} demonstrates that Self Forcing maintains strong performance across all distribution matching objectives (DMD, SiD, and GAN) and consistently outperforms alternative approaches. Two key findings emerge from these experiments: First, while baseline methods suffer significant quality degradation when shifting from chunk-wise to frame-wise architectures due to increased error accumulation (manifesting as progressive saturation or sharpness artifacts), Self Forcing maintains consistent performance across both configurations—highlighting its effectiveness at addressing exposure bias. Second, our results suggest that Teacher Forcing generally provides better initialization than Diffusion Forcing for autoregressive diffusion models. Additional qualitative comparisons in the supplementary material further illustrate these differences.
% We observe that many baseline methods exhibit clear error accumulation, . Self Forcing effectively mitigates this problem through its fundamental train-test distribution alignment that enables the model to learn from and correct its own generation errors.
% Table~\ref{tab:ablation} demonstrates that Self Forcing consistently outperforms all baselines regardless of the distribution matching objective used (DMD, SiD, or GAN). While baseline methods exhibit clear error accumulation—often manifesting as progressively increasing saturation or sharpness artifacts—Self Forcing effectively mitigates these issues by aligning train-test distributions, enabling the model to learn from and correct its own generation errors throughout the sequence. Additional qualitative comparisons in the supplementary material further illustrate this improvement in temporal consistency.
% We conduct controlled ablation studies comparing Self Forcing against alternative training approaches. Our evaluation includes: (1) conventional 50-step AR diffusion models trained with Teacher Forcing or Diffusion Forcing conditioning, and (2) few-step AR diffusion models trained with TF/DF inputs but optimized with distribution matching objectives. The latter configuration with DF+DMD precisely replicates CausVid~\cite{yin2025causvid} within our implementation framework, enabling direct comparison under identical conditions.

% Table~\ref{tab:ablation} demonstrates that Self Forcing performs consistently well across all distribution matching objectives (DMD, SiD, and GAN) and significantly outperforms all baselines on every metric. While baseline methods exhibit clear error accumulation—typically manifesting as progressively increasing saturation—Self Forcing effectively mitigates this problem through its fundamental train-test distribution alignment that enables the model to learn from and correct its own generation errors.
% We perform controlled comparisons between Self Forcing and alternative autoregressive diffusion training approaches. We evaluate: (1) AR Diffusion models trained with denoising diffusion loss using either Teacher Forcing or Diffusion Forcing conditioning, and (2) few-step AR Diffusion models trained with TF/DF inputs but using distribution matching objectives~(DMD). The latter configuration with DF + DMD essentially replicates CausVid~\cite{yin2025causvid} within our implementation framework, allowing direct comparison under identical training conditions.


% \noindent\textbf{Training efficiency.} One might expect Self Forcing training to be significantly more expensive, as it departs from the fully parallelizable training paradigm of transformers. Contrary to intuition, our experiments reveal that Self Forcing is not merely effective but actually outperforms existing strategies in training efficiency. Fig.~\ref{fig:training_time} (left) shows that Self Forcing with a single denoising step~($s=1$ in Algorithm~\ref{alg:self_forcing}) achieves lower per-iteration training time than Teacher Forcing and Diffusion Forcing.
% We attribute this result to two key factors: First, Self Forcing is not completely sequential: it processes many tokens in the same frame in parallel and we still observe high GPU utilization during training. Second, both TF and DF requires special attention masks to enforce causal dependencies, which introduces additional computational overhead even implemented with FlexAttention~\cite{Dong2024FlexAA}. We additionally show that the per-iteration cost increases gracefully with the number of denoising steps. In Fig.~\ref{fig:training_time}, we further show that Self Forcing consistently achieves superior quality under the same wall-clock training budgets compared with DF/TF, despite its sequential nature.

\begin{table}[t]
  \caption
  {
    \textbf{Ablation study.} We conduct controlled ablation studies comparing different training paradigms and distribution matching objectives under our training setup across chunk-wise~(left) and frame-wise~(right) AR models.
    Self Forcing works well with all different distribution matching objectives and consistently outperforms alternative training approaches.
  }
  \vspace{0.5em}
  \centering
  \begin{minipage}[t]{0.485\textwidth}
    \centering
    \footnotesize
    \setlength{\tabcolsep}{2pt}
    \begin{tabular}{lccc}
      \toprule
      \multirow{3}{*}{\normalsize{Chunk-wise AR}} &
      \multicolumn{3}{c}{Evaluation scores $\uparrow$}\\
    \cmidrule(lr){2-4}
                           & Total & Quality & Semantic \\
                           & Score & Score & Score \\
    \midrule
    \rowcolor{catgray}
    \multicolumn{4}{l}{\textit{Many~(50$\times$2)-step models}}\\
    Diffusion Forcing~(DF) & 82.95 & 83.66 & 80.09 \\
    Teacher Forcing~(TF)   & 83.58 & 84.34 & 80.52 \\
    \midrule
    \rowcolor{catgray}
    \multicolumn{4}{l}{\textit{Few~(4)-step models}}\\
    DF + DMD  & 82.76 & 83.49 & 79.85 \\
    % results['wan_causal_dmd_shift5_dfbaseline_ema'][500]
    TF + DMD    & 82.32 & 82.73 & 80.67 \\
    % results['wan_causal_dmd_shift5_tfbaseline_ema'][250]
    Self Forcing (Ours, DMD) & \textbf{84.31} & 85.07 & \textbf{81.28} \\
    % results['wan_causal_dmdunroll_vidprom_14bcritic_cfg3_gbeta0_warp_lrd4e_7_v3_4step_nots_shift5_fixed_ema'][600]
    Self Forcing (Ours, SiD) & 84.07 & \textbf{85.52} & 78.24 \\
    % results['wan_causal_sid_unroll_shift5_warp_mask-schedule_bs64_1500_ema'][0]
    Self Forcing (Ours, GAN) & 83.88 & 85.06 & 79.16 \\
    % results['wan_causal_r3gan-30_syn-4step_nots_ct_shift5_scratch_bs768_tsmax_ema'][100]
    \bottomrule
  \end{tabular}
  \end{minipage}
  \hfill
  \begin{minipage}[t]{0.485\textwidth}
    \centering
    \footnotesize
    \setlength{\tabcolsep}{2pt}
    \begin{tabular}{lccc}
      \toprule
      \multirow{3}{*}{\normalsize{Frame-wise AR}} &
      \multicolumn{3}{c}{Evaluation scores $\uparrow$}\\
    \cmidrule(lr){2-4}
                             & Total & Quality & Semantic \\
                             & Score & Score & Score \\
    \midrule
    \rowcolor{catgray}
    \multicolumn{4}{l}{\textit{Many~(50$\times$2)-step models}}\\
    Diffusion Forcing~(DF)   & 77.24 & 79.72 & 67.33 \\
    Teacher Forcing~(TF)     & 80.34 & 81.34 & 76.34 \\
    \midrule
    \rowcolor{catgray}
    \multicolumn{4}{l}{\textit{Few~(4)-step models}}\\
    DF + DMD                 & 80.56 & 81.02 & 78.71 \\
    % results['wan_causal_dmd_shift5_dfbaseline_block1_ema'][400]
    TF + DMD                 & 78.12 & 79.62 & 72.11 \\
    % results['wan_causal_dmd_shift5_tfbaseline_block1_ema'][500]
    Self Forcing (Ours, DMD) & \textbf{84.26} & \textbf{85.25} & \textbf{80.30} \\
    % results["wan_causal_dmdunroll_vidprom_14bcritic_cfg3_gbeta0_warp_lrd4e_7_v3_4step_nots_shift5_fixed_block1_ss_ema"][600]
    Self Forcing (Ours, SiD) & 83.54          & 84.71          & 78.86          \\
    Self Forcing (Ours, GAN) & 83.27          & 84.57          & 78.08          \\
    % results['wan_causal_r3gan-30_syn-4step_nots_ct_shift5_scratch_bs768_tsmax_block1_ema'][120]
    \bottomrule
  \end{tabular}
  \end{minipage}
  \label{tab:ablation}
  % \vspace{-1em}
\end{table}


% \subsection{Long Video Generation}


% \begin{figure}
%   \centering
%   \includegraphics[width=\textwidth,height=0.4\textwidth]{example-image-a}
%   \caption{
%     How different metrics change over time during long video generation for different models.
%   }
%   \label{fig:drift}
% \end{figure}

% \subsection{Ablation Studies}

% \begin{wraptable}{r}{0.47\textwidth}
%   \small
%   \vspace{-1.2em}
%   \centering
%   \setlength{\tabcolsep}{4.0pt}
%   \caption{
%     Comparisons for video extrapolation.
%   }
%   \label{tab:extrapolation}
%   \begin{tabular}{lcc}
%     \toprule
%     Method             & FPS   & VBench Score \\
%     \midrule
%     Recomputing KV     & 4.6   & 83.6 \\
%     Rolling KV~(naive) & 15.8  & 82.3 \\
%     Rolling KV~(ours)  & \textbf{16.1}  & \textbf{83.7} \\
%     \bottomrule
%   \end{tabular}
%   % \vspace{-1em}
% \end{wraptable}
\noindent\textbf{Rolling KV cache.} We observe that recomputing KV cache when shifting sliding window (Fig.~\ref{fig:kv}~(b)) results in significantly reduced throughput (only 4.6 FPS) when generating 10-second videos. While naive rolling KV cache maintains high throughput, it introduces severe visual artifacts, as illustrated in the examples in Appendix~\ref{app:qualitative}. By training the model to generate frames without seeing the initial image latent, we effectively mitigate these artifacts while maintaining high throughput (16.1 FPS).

% \noindent\textbf{Rolling KV cache.}
% Table~\ref{tab:extrapolation} evaluates our rolling KV cache mechanism for video extrapolation. Recomputing KV cache~(Fig.~\ref{fig:kv}~(b)) when shifting sliding window results in significantly reduced throughput~(4.6 FPS) due to redundant computation. While naive rolling KV cache maintains high throughput, it introduces visual artifacts. Though this quality drop appears modest in VBench metrics (1.4 point), qualitative examination reveals severe artifacts (see examples in Fig.~\ref{} in Appendix~\ref{app:results}). By training the model with the scenario of not seeing image latent, we can effectively mitigate this issue. Our rolling KV cache mechanism achieves a balance between efficiency and quality.
% Table~\ref{tab:extrapolation} evaluates our rolling KV cache mechanism for video extrapolation. Standard KV recomputation (Fig.~\ref{fig:kv}b) when shifting windows results in significantly reduced throughput (4.6 FPS) due to redundant computation. While naive rolling KV cache maintains high throughput, it introduces visual artifacts. Though this quality drop appears modest in VBench metrics (~1 point), qualitative examination reveals severe discontinuities (examples in supplementary material). Our attention window restriction strategy successfully mitigates these artifacts by explicitly training the model to generate coherent frames without relying on the initial image latent frame. This approach achieves an optimal balance between computational efficiency and generation quality, demonstrating the practical viability of our rolling KV cache mechanism for long-form video generation.


\noindent\textbf{Training efficiency.} One might expect Self Forcing training to be computationally prohibitive given its sequential nature that contradicts the parallelizable paradigm of transformers. Surprisingly, our experiments reveal that Self Forcing actually outperforms alternative strategies in training efficiency. As shown in Fig.~\ref{fig:training_time} (left), Self Forcing achieves comparable per-iteration training time to Teacher Forcing and Diffusion Forcing. Furthermore, Fig.~\ref{fig:training_time} (right) demonstrates that Self Forcing achieves superior quality given same wall-clock training budgets compared to both alternative approaches. Each Self Forcing experiment with DMD converges in approximately 1.5 hours on 64 H100 GPUs.

This counter-intuitive result stems from two key factors: First, while Self Forcing performs sequential rollout, it still processes all tokens within each individual frame/chunk in parallel, maintaining high GPU utilization during training. Second, TF and DF require specialized attention masking patterns to enforce causal dependencies, introducing additional computational overhead even with specialized implementations like FlexAttention~\cite{Dong2024FlexAA}.
On the other hand, Self Forcing always uses full attention during training and can leverage highly optimized attention kernels such as FlashAttention-3~\cite{shah2024flashattention}.

% \vspace{-0.5em}
\section{Discussion}
\label{sec:discussion}
In this section, we examine the broader implications of our results, discuss additional perspectives, and outline potential directions for future research.

\noindent\textbf{Fundamental limitation of the parallelizable training paradigm.} Parallelizable training has been pivotal to transformers' success by enabling efficient scaling. However, this parallelism introduces fundamental limitations.  Prior research~\cite{merrill2023parallelism} demonstrates that parallel architectures inherently limit expressiveness in sequential state-tracking problems. Our work highlights another critical limitation: parallelizable training paradigms creates misalignment between training and inference distributions, leading to the accumulation of errors over time.
We advocate a new paradigm of \textit{parallel pre-training} and \textit{sequential post-training} that combines the best of both worlds. While this paradigm shift is gaining momentum in language modeling through reinforcement learning~\cite{guo2025deepseek}, our work represents the first step towards this direction for the video domain. We believe our framework is general and can be applied to other sequence domains, especially where the data is continuous.
% \noindent\textbf{Fundamental limitation of the parallelizable training paradigm.} Parallelizable training has been pivotal to transformers' success by enabling efficient scaling. However, this parallelism introduces fundamental limitations. Prior research~\cite{merrill2023parallelism} demonstrates that parallel architectures inherently sacrifice expressiveness in sequential state-tracking problems. Our work reveals another critical limitation: parallelizable training creates inherent misalignment between training and inference distributions, leading to error accumulation during generation. We advocate a hybrid paradigm combining \textit{parallel pre-training} for efficiency with \textit{sequential post-training} for distribution alignment. While this paradigm shift is gaining momentum in language modeling through reinforcement learning~\cite{guo2025deepseek}, our work pioneers its application to video generation. The principles established here extend naturally to other sequential generation domains where training-inference distribution alignment is crucial.

\begin{figure*}[t]
  \centering
  \includegraphics[width=\textwidth]{figures/training_efficiency_combined.pdf}
  \caption{
    \textbf{Training efficiency comparison.} \textit{Left:} Per-iteration time across different chunk-wise, few-step autoregressive video diffusion training algorithms~(using DMD as the distribution matching objective). \textit{Right:} Video quality (VBench score) vs. wall clock training time.
    % \textbf{Training efficiency comparison.} \textit{Left:} per-iteration training time across chunk-wise autoregressive algorithms using DMD as the distribution matching objective. Despite its sequential nature, Self Forcing with various denoising steps ($s=1/2/3/4$ as in Algorithm~\ref{alg:self_forcing}) maintains comparable computational cost to TF/DF. \textit{Right:} Video quality (VBench score) vs. wall clock training time shows Self Forcing consistently outperforms both DF and TF given \mbox{same training budgets.}
  }
  \label{fig:training_time}
  % \vspace{-1.5em}
\end{figure*}


% Training efficiency comparison. \textit{Left:} Per-iteration training time across algorithms using DMD. Despite its sequential nature, Self Forcing with stochastically sampled denoising steps ($s$ in Algorithm~\ref{alg:self_forcing}) maintains comparable computational cost to TF/DF. \textit{Right:} Video quality (VBench score) vs. wall-clock training time shows Self Forcing consistently outperforms both Diffusion Forcing and Teacher Forcing given equal training budgets, demonstrating superior sample efficiency.

% Parallelizable training has been crucial to the success of transformer architectures, enabling efficient model scaling and accelerated training at unprecedented magnitudes. However, we argue that such parallelism comes with inherent trade-offs. Prior research~\cite{merrill2023parallelism}, for instance, demonstrates that the parallelizable nature of self-attention architectures restricts their expressive capacity, particularly in state-tracking tasks. Our work further identifies another limitation: parallelizable training paradigms often introduce discrepancies between training and inference distributions, leading to the accumulation of errors over time. To address these limitations, we advocate a hybrid approach—combining \textit{parallel pre-training} with subsequent \textit{sequential post-training}—to leverage the benefits of both methodologies. While this hybrid paradigm is becoming increasingly prevalent in natural language processing, especially with the rise of reinforcement learning and reasoning-focused models, our work is among the first to explore and adapt this direction for the video domain.

\noindent\textbf{Interplay between AR, Diffusion, and GANs.}
Autoregressive models, diffusion models, and GANs have traditionally been viewed as distinct paradigms in generative modeling. Our work highlights their complementary nature and demonstrates how they can be effectively integrated. Specifically, autoregressive and diffusion models provide complementary ways to factorize distributions~(chain-rule vs. latent-variable), which can be composed in a nested manner. The core idea behind GANs—matching the distribution of an implicit generator to the target distribution by drawing samples from the implicit generator—can be employed to train a generator powered by autoregressive-diffusion factorization.

\noindent\textbf{Limitation and future directions.}
While our method effectively mitigates error accumulation within the training context length, quality degradation remains observable when generating videos substantially longer than those seen during training. Additionally, our gradient truncation strategies—while necessary for memory efficiency—may limit the model's ability to learn long-range dependencies. Future work could explore both improved extrapolation techniques and inherently recurrent architectures like state-space models~\cite{gumamba,causalmamba} that better balance memory efficiency with long-context modeling.

\section*{Acknowledgments}
We thank Tianwei Yin, Beidi Chen, Kaiwen Zheng, Kai Zhang, Gaurav Parmar, Yi Gu, Sai Bi, and Jianming Zhang for valuable discussions. G. He and M. Zhou acknowledge the support of NSF-IIS 2212418 and NIH-R37
CA271186.
 

% TODO: cite our SSM world model paper.

% While our method significantly reduces error accumulation within training context windows, quality still degrades when generating substantially longer videos. Additionally, our gradient truncation strategies—while necessary for memory efficiency—may limit the model's capacity to learn complex temporal dependencies. Future work could explore both improved extrapolation techniques and inherently recurrent architectures like state-space models~\cite{gumamba} that better balance memory efficiency with long-range temporal modeling.

% While we primarily focused on video generation in this work, our Self Forcing algorithm can be broadly applied to other sequential domains where autoregressive diffusion hybrid models have emerged as a promising new paradigm, such as language modeling~\cite{wu2023ar,arriola2025block}, audio generation~\cite{liu2024autoregressive}, multimodal generation~\cite{zhou2025transfusion,mo2025xfusionintroducingnewmodality},
% and sequential decision making~\cite{hoeg2024streaming}.
% Future work could explore further scaling up Self Forcing to longer sequences, and developing more sophisticated memory management strategies to preserve long-range gradient flow while maintaining computational efficiency.
% The core principle of aligning training and inference distributions by having models learn from their own outputs is domain-agnostic and could benefit areas such as audio synthesis, 3D shape generation, motion prediction, and robotic trajectory planning. Any sequential generative task that suffers from exposure bias and requires efficient generation could potentially benefit from our approach. Future work could explore adapting Self Forcing to these diverse applications while developing more sophisticated gradient truncation strategies to handle even longer sequences.

% as well as other sequence domains~\cite{,,,}
{
  \small
  \bibliographystyle{plain}
  \bibliography{main}
}

\appendix

\section{Implementation Details}
\label{app:implementation}

Our implementation is largely based on the open-source code of Wan2.1~\cite{wang2025wan} and CausVid~\cite{yin2025causvid}. The attention implementation of Diffusion Forcing and Teacher Forcing baselines is based on FlexAttention~\cite{Dong2024FlexAA}, while the attention in Self Forcing is based on FlashAttention-3~\cite{shah2024flashattention}.

\paragraph{Noise schedule and model parameterization.}
Following the Wan2.1 series, we adopt the flow matching framework~\cite{lipmanflow, liu2022flow}, with time step shifting $t^\prime(k, t) = (kt / 1000) / (1 + (k-1)(t / 1000)) \cdot 1000$ and a shift factor $k=5$. The forward process is specified as $x_t = \frac{t^\prime}{1000} x + \frac{1 - t^\prime}{1000} \epsilon, \epsilon \sim \mathcal{N}(0, I)$ with $t \in [0, 1000]$.
% Following the Wan2.1 series, we adopt the flow matching framework~\cite{lipmanflow, liu2022flow}, where the forward process is specified as $x_t = \frac{t}{1000} x + \frac{1 - t}{1000} \epsilon, \epsilon \sim \mathcal{N}(0, I)$ with $t \in [0, 1000]$.
% % \TODO{Describe shifting.}
% Additionally, we employ a time step shifting $t^\prime(k) = (kt / 1000) / (1 + (k-1)(t / 1000)) \cdot 1000$.

The data prediction model is given by:
\begin{align}
    G_\theta(x, t, c) = c_{\text{skip}} \cdot \epsilon - c_{\text{out}} \cdot v_\theta(c_{\text{in}} \cdot x_t, c_{\text{noise}}(t^\prime), c).
\end{align}
We keep the preconditioning coefficients the same as the base models' configuration, i.e., $c_{\text{skip}} = c_{\text{in}} = c_{\text{out}} = 1$ and $c_{\text{noise}}(t) = t$.
Our few-step diffusion process employs a uniform 4-step schedule $[t_4, t_3, t_2, t_1]=[1000, 750, 500, 250]$.



\paragraph{Prompt preprocessing.} We use the VidProS subset from VidProM~\cite{wangvidprom}, which contains around 1M semantically unique user-written text-to-video prompts. We filter out prompts that are too short (less than 20 characters), contain command line arguments (e.g., --ar 16:9), or have a NSFW probability greater than 0.01 for any annotated category~(toxicity, obscenity, identity attack, insult, threat, and sexual explicitness). This results in a total of around 250k prompts. We then expand those prompts with \texttt{Qwen/Qwen2.5-7B-Instruct}~\cite{yang2024qwen2}, using the system prompt~(English version) provided in the open-source implementation of Wan2.1~\cite{wang2025wan}. For VBench evaluation, we similarly rewrite the test prompts using \texttt{Qwen/Qwen2.5-7B-Instruct}.
We note that the VBench results of the Wan2.1 base model are also obtained with prompt rewriting, and we report baseline results with prompt rewriting, provided that the model supports such enhancements.

% \paragraph{Distribution matching loss}
\paragraph{Training details.}
% \subsection{Training details}
Most of our training runs use 64 NVIDIA GPUs (80GB memory each) with a per-GPU batch size of 1. We implement gradient accumulation for configurations requiring a larger effective batch size than 64. Our DMD training runs take only approximately 1.5 hours to converge, while SiD/GAN training takes 2-3 hours on 64 H100 GPUs.
We initialize the real score network and critic network using the pretrained weights of the base model.
We list all other training configurations, as well as the choice of real score network and critic network for different distribution matching objectives, in Table~\ref{tab:training-hyperparameters}. We describe detailed training configurations for each distribution matching objective below.

For DMD, the gradient of the reverse Kullback-Leibler divergence is given by~\cite{yin2024one, wang2023prolificdreamer, luo2023diffinstruct}:
\begin{align}
\label{eqn:DMD-gradient}
    \nabla_\theta  \mathbb E_{t}[D_\text{KL}(p_{\theta, t} \| p_{\text{data}, t})] = -\mathbb E_{t, \hat x_t \sim q_{t|0}(\hat x_t | \hat x), \hat x \sim p_\theta(\hat x)} \left[ (s_{\text{real}}(\hat x_t, t) - s_{\text{fake}}(\hat x_t, t))\frac{\partial \hat x}{\partial \theta}\right],
    % \left(s_{\text{real}}(x_t, t)\right),
\end{align}
where $s_\text{real}(\cdot, t)$ is the score function for $p_{\text{data}, t}$, approximated by a pretrained diffusion model $f_\phi(\cdot, t)$, also referred to as the real score network, and $s_{\text{fake}}(\cdot, t)$ is the score function for $p_{\theta, t}$ and is learned through a critic network $f_\psi(\cdot, t)$ via the standard diffusion loss. The gradient in Eqn.~\eqref{eqn:DMD-gradient} is equivalent to the following loss function:
\begin{align}
    \mathcal L_{\text{DMD}}(\theta) = \mathbb E_{t, \hat x_t, \hat x}\left[ \frac{1}{2}\left\| \hat x - \mathrm{sg}\left[\hat x  - (f_\psi(\hat x_t, t) - f_\phi(\hat x_t, t))\right]\right\|^2\right],
\end{align}
where $\mathrm{sg}[\cdot]$ denotes the stop gradient operator.


Similar to the pipeline of DMD, the SiD loss is given by~\cite{zhou2024score}:
\begin{align}
    \mathcal L_{\text{SiD}}(\theta) = \mathbb E_{t, \hat x_t, \hat x}\left[ (f_\phi(\hat x_t, t) - f_\psi(\hat x_t, t))^T (f_\psi(\hat x_t, t) - \hat x) + (1 - \alpha) \| f_{\phi}(\hat x_t , t) - f_\psi(\hat x_t, t)\|^2 \right],
\end{align}
which can be shown that the case of $\alpha = 0.5$ corresponds the gradient of the Fisher divergence $\mathbb E_{t, p_{\theta, t}}[\| \nabla \log p_{\theta, t}\ - \nabla \log p_{\text{data}, t}\|^2]$ ~\cite{luo2024one, huang2024flow}. Empirically, it is observed that the second term often leads to unstable training and thus $\alpha=1$ is typically adopted for better performance~\cite{zhou2024score, zhou2025adversarial}, which is also followed in this work.

% \guande{Specify the hyperparameter corresponding to the loss}

% Please add the following required packages to your document preamble:
% \usepackage{multirow}
\begin{table}[t]
\caption{Specification of training hyperparameters}
\label{tab:training-hyperparameters}
% \guande{TODO: GAN hyperparameter}}
\vspace{0.5em}
\centering
% \setlength{\tabcolsep}{6.0pt}
\renewcommand{\arraystretch}{1.5} % Adjust this factor as needed
\resizebox{1\textwidth}{!}{
\begin{tabular}{cccc}
\toprule
Hyperparameters                         & DMD                                                                                  & SiD                                                    & GAN                                                                                 \\
\midrule
Real score network                      & Wan2.1-T2V-14B                                                                       & Wan2.1-T2V-1.3B                                        & N/A                                                                      \\
Real score CFG weight                      & 3.0                                                                       & 3.0                                        & N/A                                                                      \\
Critic network initialization           & Wan2.1-T2V-1.3B                                                                      & Wan2.1-T2V-1.3B                                        & Wan2.1-T2V-1.3B                                                                     \\
Batch size                              & 64                                                                                   & 64                                                     & 768                                                                                  \\
\multirow{2}{*}{Optimizer ($G_\theta$)} & AdamW, $\beta_1=0$, $\beta_2=0.999$,                                                                                & Adam, $\beta_1=0$, $\beta_2=0.999$,                                                    & AdamW, $\beta_1=0$, $\beta_2=0.999$,                                                                                \\
                                        & $\epsilon=1\text{e-8}$,  weight\_decay$=0.01$ & $\epsilon=1\text{e-8}$,  weight\_decay$=0$ & $\epsilon=1\text{e-8}$, weight\_decay$=0.01$ \\
\multirow{2}{*}{Optimizer ($f_\psi$)} & AdamW, $\beta_1=0$, $\beta_2=0.999$,                                                                                & Adam, $\beta_1=0$, $\beta_2=0.999$,                                                    & AdamW, $\beta_1=0$, $\beta_2=0.999$,                                                                                \\
                                        & $\epsilon=1\text{e-8}$,  weight\_decay$=0.01$ & $\epsilon=1\text{e-8}$,  weight\_decay$=0$ & $\epsilon=1\text{e-8}$, weight\_decay$=0.01$ \\
Learning rate ($G_\theta$)              & $2\text{e-6}$                                                                        & $2\text{e-6}$                                          & $2\text{e-6}$                                                                       \\
Learning rate ($f_\psi$)                & $4\text{e-7}$                                                                        & $2\text{e-6}$                                          & $2\text{e-6}$                                                                      \\
Generator/critic update ratio                & $5$                                                                        & $5$                                          & $1$                                                                      \\
EMA decay                & $0.99$                                                                        & $0.99$                                          & $0.99$                                                                      \\
\bottomrule
\end{tabular}
}
\vspace{-1em}
\end{table}

For GAN training, we add additional cross-attention layers and classification heads to the initialized critic network. We employ relativistic loss~\cite{jolicoeurrelativistic} and approximate the regularization terms~(R1 and R2) using finite difference following Seaweed-APT~\cite{lin2025diffusion}. Specifically, we perturb the noisy real/fake data with additional small Gaussian noise and encourage the discriminator output to be similar to the original one. The final training objective is defined as:
\begin{align}
    & \mathcal{L}_{\text{reg}} = \frac{1}{2} \mathbb{E}_{t, x_t, \hat{x}_t, \epsilon, \hat{\epsilon}} \left[ \| f_\psi(x_t) - f_\psi(x_t+\sigma \cdot \epsilon ) \|_2^2 + \| f_\psi(\hat{x}_t) - f_\psi(\hat{x}_t + \sigma \cdot \hat{\epsilon} ) \|_2^2 \right] \\
    & \mathcal{L}_{D}(\psi) = -\mathbb{E}_{t, x_t, \hat{x}_t} \left[ \log \left( \text{sigmoid} \left(f_\psi(x_t) - f_\psi(\hat{x}_t) \right)\right) \right] + \lambda\mathcal{L}_{\text{reg}} \\
    & \mathcal{L}_{G}(\theta) = -\mathbb{E}_{t, x_t, \hat{x}_t} \left[ \log \left( \text{sigmoid} \left(f_\psi(\hat{x}_t) - f_\psi(x_t) \right)\right) \right]
\end{align}
where $x_t\sim p_{\text{data}, t}$, $\hat{x}_t\sim p_{\theta, t}$ are the noisy real and fake data, respectively, $\epsilon$ and $\hat{\epsilon}$ are Gaussian noise sampled from $\mathcal{N}(0, 1)$, and $f_\psi$ is the critic network (discriminator) of GAN. We use $\lambda=30$, $\sigma=0.05$ for all experiments.  For a video generated from the output of the $s$-th step~(see Algorithm~\ref{alg:self_forcing} for details), we find that only sampling $t$ from $[t_{s-1}, t_{s}]$ helps stabilize the training. We also adopt a large batch size of 768 for training stability.


\section{Importance of local attention training in rolling KV cache}
\label{app:qualitative}

% Fig.~\ref{fig:comp_5s} presents qualitative comparisons of video generation results between our approach and three baselines, all sharing the same model architecture. Both Wan2.1-1.3B\cite{wang2025wan} and SkyReels-V2-1.3B~\cite{chen2025skyreels} adopt bidirectional attention and require many denoising steps, resulting in significantly slower inference compared to our method, as shown in the main paper. We also compare with CausVid~\cite{yin2025causvid}, a few-step autoregressive video diffusion model trained using the diffusion forcing (DF) paradigm. Our approach achieves competitive visual fidelity relative to the two slower models while delivering noticeably higher quality than CausVid. Specifically, videos generated by CausVid exhibit color drifting and quality degradation over time, whereas our self-forcing trained model effectively addresses this fundamental issue.




We qualitatively ablate two training settings for video extrapolation using the rolling KV cache technique. In the naive baseline, the model is trained such that every chunk always attends to the first chunk during denoising. In contrast, our proposed method restricts the attention window to prevent the model from attending to the first chunk when denoising the last chunk. As shown in Fig.~\ref{fig:comp_10s}, the naive baseline exhibits visual artifacts when extrapolating videos beyond the training context length, whereas our proposed solution mitigates this issue.



\begin{figure}[h]
  \centering
  \includegraphics[width=\textwidth]{figures/comp_10s-compressed.pdf}
  \caption{\textbf{Qualitative comparisons on video extrapolation.} We present a visual comparison between the naive baseline and our proposed technique for rolling KV cache-based video extrapolation. Compared to our method using local attention window training, extrapolated video frames from the naive baseline exhibit severe visual artifacts.}
  \label{fig:comp_10s}
\end{figure}


\section{VBench Scores Across All Dimensions}


\label{app:vbench}
\begin{wrapfigure}{r}{0.6\textwidth}
  \centering
  \includegraphics[width=\linewidth]{figures/radar.pdf}
  \caption{\textbf{VBench scores visualization.} We compare Self Forcing with SkyReels-V2~\cite{chen2025skyreels}, Wan2.1-1.3B~\cite{wang2025wan}, MAGI-1~\cite{magi1}, and CausVid~\cite{yin2025causvid} using all 16 VBench metrics.}
  \label{fig:radar}
  \vspace{-3em}
\end{wrapfigure}

In Fig.~\ref{fig:radar}, we evaluate Self Forcing~(both chunk-wise and frame-wise AR versions) using all 16 VBench metrics against representative models.
Self Forcing generally outperforms other models in terms of semantic alignment, evidenced by the high scores in scene, object class, multiple objects, and human action dimensions. Our methods also achieve good frame-wise quality, as indicated by the high scores in aesthetic quality and imaging quality.
% Our models perform slightly less well in alignment with artistic styles and camera motion, captured by appearance style and temporal style scores, respectively. This is likely because such highly specific knowledge needs to be learned during pre-training.
% , and our model is not able to perform much better than the base model in these two aspects.
Our frame-wise AR variant exhibits more dynamic motion~(high dynamic degree score) but worse temporal consistency~(worse background consistency, motion smoothness, and larger temporal flickering) than the chunk-wise AR variant.

% \begin{table*}[ht!]
%   \centering
%   \begin{adjustbox}{width=1\textwidth}
%   \begin{tabular}{lccccccccccccccccccc}
%   \toprule
%   \textbf{Method} & \textbf{Total Score} & \textbf{Quality Score} & \textbf{Semantic Score} & \textbf{Subject Consistency} & \textbf{Background Consistency} & \textbf{Temporal Flickering} & \textbf{Motion Smoothness} & \textbf{Dynamic Degree} & \textbf{Aesthetic Quality} & \textbf{Imaging Quality} & \textbf{Object Class} & \textbf{Multiple Objects} & \textbf{Human Action} & \textbf{Color} & \textbf{Spatial Relationship} & \textbf{Scene} & \textbf{Appearance Style} & \textbf{Temporal Style} & \textbf{Overall Consistency} \\
%   \midrule
%   Vchitect  & 82.24 & 83.54 & 77.06 & 96.83 & 96.66 & 98.57 & 98.98 & 63.89 & 60.41 & 65.35 & 86.61 & 68.84 & 97.20  & 87.04 & 57.55 & \textbf{56.57} & 23.73 & 25.01 & 27.57 \\
%   Jimeng    & 81.97 & 83.29 & 76.69 & 97.25 & 98.39 & 99.03 & 98.09 & 38.43 & \textbf{68.80}  & 67.09 & 89.62 & 69.08 & 90.10  & 89.05 & 77.45 & 44.94 & 22.27 & 24.7 & 27.10 \\
%   CogVideoX & 81.61 & 82.75 & 77.04 & 96.23 & 96.52 & 98.66 & 96.92 & 70.97 & 61.98 & 62.90 & 85.23 & 62.11 & 99.40  & 82.81 & 66.35 & 53.20 & \textbf{24.91} & \textbf{25.38} & \textbf{27.59} \\
%   Vidu      & 81.89 & 83.85 & 74.04 & 94.63 & 96.55 & 99.08 & 97.71 & 82.64 & 60.87 & 63.32 & 88.43 & 61.68 & 97.40  & 83.24 & 66.18 & 46.07 & 21.54 & 23.79 & 26.47 \\
%   Kling     & 81.85 & 83.39 & 75.68 & \textbf{98.33} & 97.60 & \textbf{99.30} & \textbf{99.40} & 46.94 & 61.21 & 65.62 & 87.24 & 68.05 & 93.40  & \textbf{89.90}  & 73.03 & 50.86 & 19.62 & 24.17 & 26.42 \\
%   CogVideoX1.5-5B & 82.17 & 82.78 & \textbf{79.76} & 96.87 & \textbf{97.35} & 98.88 & 98.31 & 50.93 & 62.79 & 65.02 & 87.47 & 69.65 & 97.20  & 87.55 & \textbf{80.25} & 52.91 & 24.89 & 25.19 & 27.30 \\
%   Gen-3     & 82.32 & 84.11 & 75.17 & 97.10  & 96.62 & 98.61 & 99.23 & 60.14 & 63.34 & 66.82 & 87.81 & 53.64 & 96.40  & 80.90  & 65.09 & 54.57 & 24.31 & 24.71 & 26.69 \\
%   \textbf{Self Forcing}      & \textbf{84.27} & \textbf{85.65} & 78.75  & 97.53 & 97.19 & 96.24 & 98.05 & \textbf{92.69}  & 64.15 & \textbf{68.88}  & \textbf{92.99}  & \textbf{72.15} & \textbf{99.80}  & 80.17 & 64.65 & 56.58 & 24.27 & 25.33 & 27.51 \\
%   % \midrule
%   % Bidirectional Teacher & 82.10 & 83.12 & 78.01 & 90.26 & 95.07 & 95.16 &
%   % 97.88 & 92.22 & 65.18 & 66.54 & 90.30 & 71.62 & 99.80 & 79.50 & 64.20 & 53.11& 24.19 & 26.03 & 27.63 \\
%   \bottomrule
%   \end{tabular}
%   \end{adjustbox}
%   \caption{Full comparison on VBench using all 16 metrics. The best scores are in bold. Please zoom in for details. TODO: Update numbers}
%   \label{tab:vbench_long}
% \end{table*}




\section{Broader Societal Impact}
\label{app:broader_impact}
Generative modeling—particularly for videos—carries significant potential for misuse. It can lead to serious societal consequences, most notably the spread of disinformation through deepfakes that become increasingly difficult to distinguish from authentic content. Additionally, these models can reinforce harmful stereotypes and amplify existing societal biases without careful governance and responsible deployment.

Our research on real-time video generation creates additional complexities, as it removes one of the practical barriers (computational cost) that currently limits widespread misuse. While our methods enable positive applications like creative content production and accessibility tools, we acknowledge the dual-use nature of this technology and encourage continued research into detection methods, watermarking techniques, and policy frameworks that can help mitigate potential harms.



\section{User Study Details}
\label{app:user_study}
% We conduct user studies to evaluate different methods based on user preferences on Amazon MTurk.
In the user preference study, we show users two videos side by side using the same text prompt.
We ask the users to select the one that is overall better, considering both quality and prompt alignment. Detailed instructions are shown in Fig.~\ref{fig:user_study_instruction}.
We use all 1003 prompts from MovieGenBench~\cite{polyak2024movie} and each prompt is evaluated by a single user.
% We provide a compensation of \$0.03 per HIT.


\begin{figure}[h!]
  \centering
  \includegraphics[width=\textwidth]{figures/user_study.jpg}
  \caption{\textbf{User study instruction screenshots.}}
  \label{fig:user_study_instruction}
\end{figure}
% \input{checklist.tex}


\end{document}